\noindent\textbf{Operaciones y resultados intermedios.}

\begin{enumerate}

\item Área equivalente que soporta tensión axial.

\[\fitmath{y_{1}=2\cdot 65+2\cdot 75+2\cdot 110\approx 500\,\mathrm{mm^{2}}}\]

\item Propiedades geométricas de la sección (centroide y momentos de inercia).

\[\fitmath{z_G=2.8\,\mathrm{mm},\quad y_G=24.2\,\mathrm{mm}}\]

\item Momentos de inercia del modelo en los ejes centroidales.

\[\fitmath{I_z=594513.333\,\mathrm{mm}^{4},\quad I_y=460413.333\,\mathrm{mm}^{4},\quad I_{yz}=-33880\,\mathrm{mm}^{4}}\]

\item Centro de cortadura a partir del equilibrio de los flujos; coordenadas del modelo.

\[\fitmath{z_C=0\,\mathrm{mm},\quad y_C=0\,\mathrm{mm}}\]

\item Tensión máxima |q|/t en pared alma.

\[\fitmath{y_{5}=\frac{\left|\frac{8133213402}{142874479}\right|}{2}\approx 28.46279\,\mathrm{MPa}}\]

\end{enumerate}
